\allowdisplaybreaks
\begin{center}
	\hrule
	\vspace{.4cm}
 \begin{tabular*}{\textwidth}{@{}l@{}|@{\extracolsep{0.6in}}r@{}}%
\parbox{4.25in}{\raggedright{\includegraphics[width=.9\linewidth]{ictp-pwf.pdf}}} &
\parbox[c][]{4in}{{\Large\textbf{Md Akiful Islam Zawad} \par}
                    { Brac University \par}
                    { KHA 224, Progati Sarani, Merul Badda, \par}
                    { Dhaka 1212, Bangladesh \par}
                    { \href{ext.akiful.islam@bracu.ac.bd}{ext.akiful.islam@bracu.ac.bd}} \par}
\end{tabular*}\vspace{.3in}
\definecolor{ceruleanblue}{rgb}{0.16, 0.32, 0.6}
	\LARGE\scshape\textbf{\textcolor{ceruleanblue}{Physics for Bangladesh QFT School}}
\end{center}
\hrule\vspace{.25in}
{\large\textbf{Problem Sheet:}\ \textsc{6} \hspace{\hfill} \large\textbf{Due Date:} \today\\
	\hrule}
 %----------------------------
\paragraph*{Problem Set 2} %\hfill \newline
\\
Electroweak theory. Consider a toy model of the Weinberg-Salam theory with gauge group $\text{SU(2)} \otimes \text{U(1)}$, a complex scalar doublet, and two Dirac fermions.
\begin{align}
    L = - \frac{1}{4} \mathbf{F}_{\mu\nu} \cdot \mathbf{F}^{\mu\nu} - \frac{1}{4} B_{\mu\nu}B^{\mu\nu} + (D_\mu \varphi)^\dagger \cdot D^\mu \varphi + i \bar{\psi} \gamma^\mu D_\mu \psi - \frac{\lambda}{4} \left( \varphi^\dagger \cdot \varphi - \frac{v^2}{2} \right)^2 - \left( \Gamma_2 \bar{\psi} \varphi P_R \psi_2 + \Gamma_1 \bar{\psi} \varphi^c P_R \psi_1 + \text{h.c.} \right), 
\end{align}
where:
\begin{align}
    \mathbf{F}_{\mu\nu} &= \partial_\mu \mathbf{W}_\nu - \partial_\nu \mathbf{W}_\mu - g \mathbf{W}_\mu \times \mathbf{W}_\nu , \quad B_{\mu\nu} = \partial_\mu B_\nu - \partial_\nu B_\mu, \\
    \varphi &= \begin{pmatrix} \varphi_1 \\ \varphi_2 \end{pmatrix}, \quad \psi = \begin{pmatrix} \psi_1 \\ \psi_2 \end{pmatrix}, \\
    \varphi^c &= i\sigma_2 \varphi^*, \quad P_{L,R} = \frac{1 \mp \gamma^5}{2}, \\
    D_\mu \varphi &= \partial_\mu \varphi + i \left( g \mathbf{W}_\mu \cdot \frac{\sigma}{2} + g^\prime B_\mu Y \right) \varphi, \\
    D_\mu \psi &= \partial_\mu \psi + i \left( g \mathbf{W}_\mu \cdot \frac{\sigma}{2} + g^\prime B_\mu y \right) P_L \psi + ig^\prime B_\mu (y + Y \sigma_3) P_R \psi,
\end{align}
Here, boldface quantities indicate triplets, for instance:
\begin{align}
    \mathbf{W}_\mu = \begin{pmatrix} W_\mu^1 \\\\ W_\mu^2 \\\\ W_\mu^3 \end{pmatrix},
\end{align}
and $\times$ is the standard vector product while $\cdot$ indicates the scalar product. The scalar and fermions are written as a 2x2 vector, which indicates their SU(2) transformations (and so are subject to the action of, e.g., Pauli matrices $\sigma$). In addition, $\lambda, v, Y, y, \Gamma_1$ and $\Gamma_2$ are positive numbers (Higgs couplings, hypercharges, and Yukawa couplings).

Show that this Lagrangian density is gauge invariant and rule out the existence of bare mass terms on the grounds of gauge invariance. Then, show that, upon SSB, mass terms are produced for some of the fields. Identify these fields and compute their masses. In order to have standard kinetic terms, you will have to perform field redefinitions. For instance, you should find:
\begin{align}
    m_W^2 W^{+\dagger}_\mu W^\mu + \frac{1}{2} m_Z^2 Z_\mu Z^\mu,
\end{align}
where:
\begin{align}
    W_\mu^\pm = \frac{1}{\sqrt{2}} (W_\mu^1 \mp i W_\mu^2), \quad Z_\mu = \cos(\theta_W) W_\mu^3 - \sin(\theta_W) B_\mu,
\end{align}
and the Weinberg angle is defined as:
\begin{align}
    \cos(\theta_W) = \frac{g}{\sqrt{g^2 + g^\prime ^2}}, \quad \sin(\theta_W) = \frac{g^\prime }{\sqrt{g^2 + g^\prime ^2}}.
\end{align}

How is $\displaystyle\frac{m_Z}{m_W}$ related to $\theta_W$?
\bigskip\bigskip\hline\hline\bigskip
\subsection*{Solution}
Under the gauge group $\text{SU(2)} \otimes \text{U(1)}$, the fields and covariant derivatives transform as follows:

For the scalar field doublet $\varphi$:
\begin{align}
    \varphi &\to e^{\displaystyle i \alpha(x) Y} e^{\displaystyle i \boldsymbol{\theta}(x) \cdot \frac{\boldsymbol{\sigma}}{2}} \varphi
\end{align}
where $Y$ is the hypercharge and $\boldsymbol{\sigma}$ are the Pauli matrices.

For the fermions $\psi$:
\begin{align}
    \psi &\to e^{\displaystyle i \alpha(x) y} e^{\displaystyle i \boldsymbol{\theta}(x) \cdot \frac{\boldsymbol{\sigma}}{2}} \psi
\end{align}
where $y$ is the fermion hypercharge.

For the gauge fields $W_\mu$ (for SU(2)) and $B_\mu$ (for U(1)):
\begin{align}
    \mathbf{W}_\mu &\to \mathbf{W}_\mu^\prime = \mathbf{W}_\mu + \frac{1}{g} \partial_\mu \boldsymbol{\theta} + \boldsymbol{\theta} \times \mathbf{W}_\mu \\
    B_\mu &\to B_\mu^\prime = B_\mu + \frac{1}{g^\prime } \partial_\mu \alpha
\end{align}
The field strength tensors transform covariantly:
\begin{align}
    \mathbf{F}_{\mu\nu} &\to \mathbf{F}_{\mu\nu}^\prime = e^{\displaystyle i \boldsymbol{\theta}(x) \cdot \frac{\boldsymbol{\sigma}}{2}} \mathbf{F}_{\mu\nu} e^{\displaystyle -i \boldsymbol{\theta}(x) \cdot \frac{\boldsymbol{\sigma}}{2}} \\
    B_{\mu\nu} &\to B_{\mu\nu}^\prime = B_{\mu\nu}
\end{align}
We now study the terms in the Lagrangian:
\begin{enumerate}
    \item [(i)] \textbf{Gauge Kinetic Terms for the Gauge Fields}

    The terms involving the field strength tensors $\mathbf{F}_{\mu\nu}$ and $B_{\mu\nu}$ are:
    \begin{align}
        L_{\text{gauge}} = - \frac{1}{4} \mathbf{F}_{\mu\nu} \cdot \mathbf{F}^{\mu\nu} - \frac{1}{4} B_{\mu\nu} B^{\mu\nu}\label{eq:L-gauge-def}
    \end{align}
    The field strength tensor $\mathbf{F}_{\mu\nu}$ for the $\text{SU(2)}$ gauge field $W_\mu$ is given by:
    \begin{align}
        \mathbf{F}_{\mu\nu} &= \partial_\mu \mathbf{W}_\nu - \partial_\nu \mathbf{W}_\mu - g \mathbf{W}_\mu \times \mathbf{W}_\nu 
    \end{align}
    where $W_\mu$ is a triplet of gauge fields represented in column form:
    \begin{align*}
        \mathbf{W}_\mu = \begin{pmatrix} W_\mu^1 \\\\ W_\mu^2 \\\\ W_\mu^3 \end{pmatrix}
    \end{align*}
    Here, the cross product $\times$ denotes the non-Abelian interaction among the SU(2) gauge fields. Under a gauge transformation with parameter $\boldsymbol{\theta}(x)$ for $\text{SU(2)}$, the gauge field transforms as:
    \begin{align*}
        \mathbf{W}_\mu &\to \mathbf{W}_\mu^\prime = \mathbf{W}_\mu + \frac{1}{g} \partial_\mu \boldsymbol{\theta} + \boldsymbol{\theta} \times W_\mu
    \end{align*}
    where $\boldsymbol{\theta}(x)$ is a triplet corresponding to the gauge transformation.

    Substituting this into the expression for $\mathbf{F}_{\mu\nu}$ (\ref{eq:L-gauge-def}), we calculate how the field strength tensor transforms:
    \begin{align}
        \mathbf{F}_{\mu\nu}^\prime &= \partial_\mu \mathbf{W}_\nu ^\prime - \partial_\nu \mathbf{W}_\mu^\prime - g \mathbf{W}_\mu^\prime \times \mathbf{W}_\nu ^\prime \notag \\
        &= \partial_\mu \left( \mathbf{W}_\nu + \frac{1}{g} \partial_\nu \boldsymbol{\theta} + \boldsymbol{\theta}  \mathbf{W}_\nu \right)
        - \partial_\nu \left( \mathbf{W}_\mu + \frac{1}{g} \partial_\mu \boldsymbol{\theta} + \boldsymbol{\theta} \times \mathbf{W}_\mu \right) - g \left( \mathbf{W}_\mu + \frac{1}{g} \partial_\mu \boldsymbol{\theta} + \boldsymbol{\theta} \times \mathbf{W}_\mu \right) \times \left( \mathbf{W}_\nu + \frac{1}{g} \partial_\nu \boldsymbol{\theta} + \boldsymbol{\theta} \times \mathbf{W}_\nu \right) \notag \\
        &= \left( \partial_\mu \mathbf{W}_\nu - \partial_\nu \mathbf{W}_\mu - g \mathbf{W}_\mu \times \mathbf{W}_\nu \right)
        + \boldsymbol{\theta} \times \mathbf{F}_{\mu\nu} \notag \\[5pt]
        &= \mathbf{F}_{\mu\nu} + \boldsymbol{\theta} \times \mathbf{F}_{\mu\nu}\label{eq:F-mu-nu-transformation}
    \end{align}
    The gauge field $W_\mu$ is represented in terms of Pauli matrices as:
    \begin{align}
        \mathbf{W}_\mu = W_\mu^a \frac{\sigma^a}{2}
    \end{align}
    where $a = 1, 2, 3$ labels the components of the SU(2) triplet, and $\sigma^a$ are the Pauli matrices.
    
    Similarly, the field strength tensor $\mathbf{F}_{\mu\nu}$ is written as:
    \begin{align}
        \mathbf{F}_{\mu\nu} = \mathbf{F}_{\mu\nu}^a \frac{\sigma^a}{2}
    \end{align}
    We start with the covariant transformation of the field strength tensor $\mathbf{F}_{\mu\nu}$ under a gauge transformation:
    \begin{align}
        \mathbf{F}_{\mu\nu}^\prime = e^{\displaystyle \displaystyle i \boldsymbol{\theta}(x) \cdot \frac{\boldsymbol{\sigma}}{2}} \mathbf{F}_{\mu\nu} e^{\displaystyle -i \boldsymbol{\theta}(x) \cdot \frac{\boldsymbol{\sigma}}{2}}
    \end{align}    
    We will use the Baker-Campbell-Hausdorff formula to expand this transformation. The Baker-Campbell-Hausdorff formula is:
    \begin{align*}
        e^A e^B = \exp\left( A + B + \frac{1}{2} [A, B] + \frac{1}{12} [A, [A, B]] - \frac{1}{12} [B, [A, B]] + \dots \right)
    \end{align*}
    For an infinitesimal gauge transformation where the gauge parameter $\boldsymbol{\theta}(x)$ is small, we approximate the transformation using the first few terms of the expansion. Let:
    \begin{align*}
        A = i \boldsymbol{\theta}(x) \cdot \frac{\boldsymbol{\sigma}}{2}, \quad B = -i \boldsymbol{\theta}(x) \cdot \frac{\boldsymbol{\sigma}}{2}
    \end{align*}
    Thus, the transformation becomes:
    \begin{align*}
        \mathbf{F}_{\mu\nu}^\prime &= e^{\displaystyle\left( i \boldsymbol{\theta}(x) \cdot \frac{\boldsymbol{\sigma}}{2} \right)} \mathbf{F}_{\mu\nu} e^{\displaystyle\left( -i \boldsymbol{\theta}(x) \cdot \frac{\boldsymbol{\sigma}}{2} \right)} \notag \\
        &= \left( 1 + i \boldsymbol{\theta}(x) \cdot \frac{\boldsymbol{\sigma}}{2} + \dots \right) \mathbf{F}_{\mu\nu} \left( 1 - i \boldsymbol{\theta}(x) \cdot \frac{\boldsymbol{\sigma}}{2} + \dots \right)\notag\\
        &= \mathbf{F}_{\mu\nu} + i \left( \boldsymbol{\theta}(x) \cdot \frac{\boldsymbol{\sigma}}{2} \right) \mathbf{F}_{\mu\nu} - i \mathbf{F}_{\mu\nu} \left( \boldsymbol{\theta}(x) \cdot \frac{\boldsymbol{\sigma}}{2} \right) \notag \\
    &= \mathbf{F}_{\mu\nu} + \left( \boldsymbol{\theta}(x) \times \mathbf{F}_{\mu\nu} \right) \cdot \frac{\boldsymbol{\sigma}}{2}
\end{align*}
This confirms the transformation in (\ref{eq:F-mu-nu-transformation}) and implies that the gauge kinetic term involving $\mathbf{F}_{\mu\nu}$: $\displaystyle-\frac{1}{4} \mathbf{F}_{\mu\nu} \cdot \mathbf{F}^{\mu\nu}$ is invariant under $\text{SU(2)}$ gauge transformations because $\mathbf{F}_{\mu\nu}$ transforms covariantly, and the scalar product $\cdot$ is invariant under such transformations.

Next, we see the field strength tensor for the U(1) gauge field $B_\mu$ is:
\begin{align}
    B_{\mu\nu} &= \partial_\mu B_\nu - \partial_\nu B_\mu\label{eq:B-mu-nu-def}
\end{align}
Under a U(1) gauge transformation, the gauge field $B_\mu$ transforms as:
\begin{align}
    B_\mu &\to B_\mu^\prime = B_\mu + \frac{1}{g^\prime } \partial_\mu \alpha\label{eq:B-mu-nu-transformation}
\end{align}
where $\alpha(x)$ is the U(1) gauge transformation parameter.

Substituting this into the expression for $B_{\mu\nu}$ (\ref{eq:B-mu-nu-def}), we find that:
\begin{align}
    B_{\mu\nu}^\prime &= \partial_\mu B_\nu^\prime - \partial_\nu B_\mu^\prime \notag \\
    &= \partial_\mu \left( B_\nu + \frac{1}{g^\prime } \partial_\nu \alpha \right) - \partial_\nu \left( B_\mu + \frac{1}{g^\prime } \partial_\mu \alpha \right) \notag \\
    &= \partial_\mu B_\nu - \partial_\nu B_\mu \notag \\[3pt]
    &= B_{\mu\nu}
\end{align}

Thus, $B_{\mu\nu}$ is invariant under the U(1) gauge transformation.

Therefore, the gauge kinetic term for U(1): $\displaystyle- \frac{1}{4} B_{\mu\nu} B^{\mu\nu}$ is gauge invariant.

\item [(ii)] \textbf{Kinetic Terms for the Scalar Field}

The kinetic term for the scalar field $\varphi$ is:
\begin{align}
    L_{\varphi} = (D_\mu \varphi)^\dagger D^\mu \varphi
\end{align}
where the covariant derivative is defined as:
\begin{align}
    D_\mu \varphi = \partial_\mu \varphi + i \left( g \mathbf{W}_\mu \cdot \frac{\sigma}{2} + g^\prime B_\mu Y \right) \varphi
\end{align}
Here, $W_\mu$ is the triplet gauge field for $\text{SU(2)}$, $B_\mu$ is the $\text{U(1)}$ gauge field, $\sigma^a$ are the Pauli matrices, $g$ and $g^\prime $ are the coupling constants for $\text{SU(2)}$ and $\text{U(1)}$, and $Y$ is the hypercharge.

Under a local $\text{SU(2)} \otimes \text{U(1)}$ gauge transformation with parameter $\boldsymbol{\theta}(x)$ for $\text{SU(2)}$ and $\alpha(x)$ for $\text{U(1)}$, the scalar field $\varphi$ transforms covariantly:
\begin{align}
    \varphi(x) \to \varphi^\prime (x) &= e^{\displaystyle i \alpha(x) Y} e^{\displaystyle i \boldsymbol{\theta}(x) \cdot \frac{\boldsymbol{\sigma}}{2}} \varphi(x)
\end{align}
Now we show that the covariant derivative $D_\mu \varphi$ also transforms covariantly under the same gauge transformation. The covariant derivative is:
\begin{align}
    D_\mu \varphi = \partial_\mu \varphi + i \left( g \mathbf{W}_\mu \cdot \frac{\sigma}{2} + g^\prime B_\mu Y \right) \varphi\label{eq:covariant-derivative-def}
\end{align}
Under a gauge transformation, the gauge fields $W_\mu$ and $B_\mu$ transform as:
\begin{align}
    \mathbf{W}_\mu \to \mathbf{W}_\mu^\prime &= e^{\displaystyle i \boldsymbol{\theta}(x) \cdot \frac{\sigma}{2}} \left( \mathbf{W}_\mu - \frac{1}{g} \partial_\mu \boldsymbol{\theta}(x) \right) e^{\displaystyle -i \boldsymbol{\theta}(x) \cdot \frac{\sigma}{2}} \\
    B_\mu \to B_\mu^\prime &= B_\mu - \frac{1}{g^\prime } \partial_\mu \alpha(x)
\end{align}
Let’s apply the transformation to the covariant derivative:
\begin{align*}
    D_\mu^\prime \varphi^\prime &= \partial_\mu \left( e^{\displaystyle i \alpha(x) Y} e^{\displaystyle i \boldsymbol{\theta}(x) \cdot \frac{\boldsymbol{\sigma}}{2}} \varphi \right) + i \left( g \mathbf{W}_\mu^\prime \cdot \frac{\sigma}{2} + g^\prime B_\mu^\prime Y \right) \left( e^{\displaystyle i \alpha(x) Y} e^{\displaystyle i \boldsymbol{\theta}(x) \cdot \frac{\boldsymbol{\sigma}}{2}} \varphi \right)\notag\\
    &= \left( \partial_\mu e^{\displaystyle i \alpha(x) Y} \right) e^{\displaystyle i \boldsymbol{\theta}(x) \cdot \frac{\boldsymbol{\sigma}}{2}} \varphi + e^{\displaystyle i \alpha(x) Y} \left( \partial_\mu e^{\displaystyle i \boldsymbol{\theta}(x) \cdot \frac{\boldsymbol{\sigma}}{2}} \right) \varphi + e^{\displaystyle i \alpha(x) Y} e^{\displaystyle i \boldsymbol{\theta}(x) \cdot \frac{\boldsymbol{\sigma}}{2}} \partial_\mu \varphi\notag\\
    &\qquad~~~~+ i \left( g \mathbf{W}_\mu^\prime \cdot \frac{\sigma}{2} + g^\prime B_\mu^\prime Y \right) e^{\displaystyle i \alpha(x) Y} e^{\displaystyle i \boldsymbol{\theta}(x) \cdot \frac{\boldsymbol{\sigma}}{2}} \varphi
\end{align*}
Using the fact that the gauge parameters $\alpha(x)$ and $\boldsymbol{\theta}(x)$ are infinitesimal, we can simplify the terms involving the derivatives of the exponential factors. The transformed covariant derivative becomes:
\begin{align*}
    D_\mu^\prime \varphi^\prime &= e^{\displaystyle i \alpha(x) Y} e^{\displaystyle i \boldsymbol{\theta}(x) \cdot \frac{\boldsymbol{\sigma}}{2}} \left( \partial_\mu \varphi + i \left( g \mathbf{W}_\mu \cdot \frac{\sigma}{2} + g^\prime B_\mu Y \right) \varphi \right)\label{eq:covariant-transforms-varphi}
\end{align*}
Thus, we find that the covariant derivative transforms covariantly (\ref{eq:covariant-derivative-def}):
\begin{align}
    D_\mu \varphi \to D_\mu^\prime \varphi^\prime &= e^{\displaystyle i \alpha(x) Y} e^{\displaystyle i \boldsymbol{\theta}(x) \cdot \frac{\boldsymbol{\sigma}}{2}} D_\mu \varphi
\end{align}
Using the fact that $D_\mu \varphi$ transforms covariantly, we now show that the kinetic term is gauge-invariant. Under a gauge transformation, we have:
\begin{align}
    (D_\mu \varphi)^\prime &= (D_\mu \varphi) \to e^{\displaystyle i \alpha(x) Y} e^{\displaystyle i \boldsymbol{\theta}(x) \cdot \frac{\boldsymbol{\sigma}}{2}} D_\mu \varphi
\end{align}

For the Hermitian conjugate $(D_\mu \varphi)^\dagger$, we have:
\begin{align}
    (D_\mu \varphi)^\dagger \to (D_\mu \varphi)^\dagger e^{\displaystyle -i \alpha(x) Y} e^{\displaystyle -i \boldsymbol{\theta}(x) \cdot \frac{\boldsymbol{\sigma}}{2}}
\end{align}
Now, the transformed kinetic term becomes:
\begin{align*}
    (D_\mu^\prime \varphi^\prime )^\dagger D^\mu^\prime \varphi^\prime &= \left( (D_\mu \varphi)^\dagger e^{\displaystyle -i \alpha(x) Y} e^{\displaystyle-i \boldsymbol{\theta}(x) \cdot \frac{\boldsymbol{\sigma}}{2}} \right) \left( e^{\displaystyle i \alpha(x) Y} e^{\displaystyle i \boldsymbol{\theta}(x) \cdot \frac{\boldsymbol{\sigma}}{2}} D^\mu \varphi \right)
\end{align*}
Since $\displaystyle e^{\displaystyle i \alpha(x) Y} e^{\displaystyle i \boldsymbol{\theta}(x) \cdot \frac{\boldsymbol{\sigma}}{2}}$ and its inverse cancel each other, the kinetic term simplifies to:
\begin{align*}
    (D_\mu^\prime \varphi^\prime )^\dagger D^\mu^\prime \varphi^\prime &= (D_\mu \varphi)^\dagger D^\mu \varphi
\end{align*}
Therefore, the kinetic terms for the scalar field are gauge invariant.
\item [(iii)] \textbf{Kinetic terms for the fermions}

The kinetic term for the fermions is:
\begin{align}
    L_{\psi} = i \bar{\psi} \gamma^\mu D_\mu \psi,\label{eq:L-psi-def}
\end{align}
where the covariant derivative $D_\mu \psi$ is:
\begin{align}
    D_\mu \psi = \partial_\mu \psi + i \left( g \mathbf{W}_\mu \cdot \frac{\sigma}{2} + g^\prime B_\mu y \right) P_L \psi + i g^\prime B_\mu (y + Y \sigma_3) P_R \psi
\end{align}
In the Standard Model, the fermion fields $\psi$ transform under the $\text{SU(2)} \otimes \text{U(1)}_Y$ gauge group. The left-handed and right-handed components of $\psi$ transform differently under this gauge symmetry.

The \textbf{left-handed fermions} transform as a doublet under $\text{SU(2)}$ and have hypercharge $y$ under $\text{U(1)}_Y$:
\begin{align}
    \psi_L(x) \to U(x) \psi_L(x) = e^{\displaystyle i \alpha_a(x) \frac{\sigma^a}{2} + i \beta(x) y} \psi_L(x)
\end{align}
The \textbf{right-handed fermions} transform as singlets under $\text{SU(2)}$ and have hypercharge $y + Y \sigma_3$:
\begin{align}
    \psi_R(x) \to e^{\displaystyle i \beta(x) (y + Y \sigma_3)} \psi_R(x)
\end{align}
The covariant derivative acting on the fermion fields $\psi$ is given by:
\begin{align}
    D_\mu \psi &= \partial_\mu \psi + i \left( g \mathbf{W}_\mu \cdot \frac{\sigma}{2} + g^\prime B_\mu y \right) P_L \psi + i g^\prime B_\mu (y + Y \sigma_3) P_R \psi
\end{align}
The term $\displaystyle g \mathbf{W}_\mu \cdot \frac{\sigma}{2}$ represents the coupling of the $\text{SU(2)}$ gauge bosons $W_\mu^a$ to the left-handed fermions, with $g$ as the coupling constant and $\displaystyle\frac{\sigma^a}{2}$ as the generators of $\text{SU(2)}$.

The terms involving $g^\prime $ represent the coupling of the $\text{U(1)}_Y$ gauge boson $B_\mu$ to both left-handed and right-handed fermions. The left-handed fermions couple with hypercharge $y$, while the right-handed fermions couple with hypercharge $y + Y \sigma_3$.

Under a gauge transformation, the covariant derivative must transform covariantly, i.e.:
\begin{align}
    D_\mu \psi \to U(x) D_\mu \psi
\end{align}
where $U(x)$ is the local gauge transformation matrix:
\begin{align}
    U(x) = e^{\displaystyle i \alpha_a(x) \frac{\sigma^a}{2} + i \beta(x) Y}
\end{align}
To check the gauge invariance of the covariant derivative, we examine how it transforms under a gauge transformation.
\begin{itemize}
    \item The left-handed component $\psi_L(x)$ transforms as:
   \begin{align}
       \psi_L(x) \to e^{\displaystyle i \alpha_a(x) \frac{\sigma^a}{2} + i \beta(x) y} \psi_L(x)
   \end{align}
   The covariant derivative acting on $\psi_L(x)$ becomes:
   \begin{align}
       D_\mu \psi_L(x) &= \partial_\mu \psi_L(x) + i \left( g W_\mu^a \frac{\sigma^a}{2} + g^\prime  B_\mu y \right) \psi_L(x)
   \end{align}
   Under a gauge transformation, the gauge fields $W_\mu^a$ and $B_\mu$ transform as:
   \begin{align}
       W_\mu^a &\to W_\mu^a + \frac{1}{g} \partial_\mu \alpha^a + \epsilon^{abc} \alpha^b W_\mu^c \\
       B_\mu &\to B_\mu + \frac{1}{g^\prime } \partial_\mu \beta
   \end{align}
   The covariant derivative transforms covariantly under this gauge transformation:
   \begin{align}
       D_\mu \psi_L(x) \to e^{\displaystyle i \alpha_a(x) \frac{\sigma^a}{2} + i \beta(x) y} D_\mu \psi_L(x)
   \end{align}
\item The right-handed component $\psi_R(x)$ transforms as:
   \begin{align}
       \psi_R(x) \to e^{\displaystyle i \beta(x) (y + Y \sigma_3)} \psi_R(x)
   \end{align}
   The covariant derivative acting on $\psi_R(x)$ becomes:
   \begin{align}
       D_\mu \psi_R(x) &= \partial_\mu \psi_R(x) + i g^\prime  B_\mu (y + Y \sigma_3) \psi_R(x)
   \end{align}
   Under a gauge transformation, the covariant derivative transforms as:
   \begin{align}
       D_\mu \psi_R(x) \to e^{\displaystyle i \beta(x) (y + Y \sigma_3)} D_\mu \psi_R(x)
   \end{align}
\end{itemize}
The fermion kinetic term is:
\begin{align}
    L_{\psi} = i \bar{\psi} \gamma^\mu D_\mu \psi
\end{align}

The fermion fields $\psi$ transform covariantly under the gauge group, meaning:
\begin{align}
    \psi(x) \to U(x) \psi(x), \quad \bar{\psi}(x) \to \bar{\psi}(x) U^\dagger(x)
\end{align}

Since both the fermion fields and the covariant derivative transform covariantly, the fermion kinetic term is gauge-invariant:
\begin{align}
    L_{\psi} &\to i \bar{\psi} U^\dagger(x) \gamma^\mu D_\mu (U(x) \psi) \notag \\
    &= i \bar{\psi} \gamma^\mu D_\mu \psi
\end{align}
This shows that the fermion kinetic term remains invariant under the full $\text{SU(2)} \otimes \text{U(1)}_Y$ gauge transformations.



This term is gauge-invariant because both the covariant derivative and the fermion fields transform covariantly.
\item [(iv)] \textbf{Scalar Potential Term}:

The scalar potential term is:
\begin{align}
    L_{\text{potential}} = - \frac{\lambda}{4} \left( \varphi^\dagger \varphi - \frac{v^2}{2} \right)^2\label{eq:L-scalar-potential}
\end{align}
We start by applying a local $U(1)$ gauge transformation. The scalar field $\varphi(x)$ transforms as:
\begin{align}
    \varphi(x) &\to e^{\displaystyle i \alpha(x)} \varphi(x)\label{eq:phi-potential-transformation}
\end{align}
where $\alpha(x)$ is a real, space-time-dependent parameter representing the local gauge transformation.

Similarly, the complex conjugate of the scalar field $\varphi^\dagger(x)$ (\ref{eq:phi-potential-transformation}) transforms as:
\begin{align}
    \varphi^\dagger(x) &\to e^{-i \alpha(x)} \varphi^\dagger(x)
\end{align}
Next, we examine how the combination $\varphi^\dagger \varphi$ transforms. Under the gauge transformation:
\begin{align*}
    \varphi^\dagger(x) \varphi(x) &\to e^{\displaystyle -i \alpha(x)} \varphi^\dagger(x) e^{\displaystyle i \alpha(x)} \varphi(x)
\end{align*}
Since the exponentials $e^{\displaystyle i \alpha(x)}$ and $e^{\displaystyle -i \alpha(x)}$ cancel each other, we are left with:
\begin{align*}
    \varphi^\dagger(x) \varphi(x) &\to \varphi^\dagger(x) \varphi(x)
\end{align*}
This shows that $\varphi^\dagger \varphi$ is invariant under the local gauge transformation. Now, we analyze the scalar potential term. It is given by:
\begin{align*}
    \left( \varphi^\dagger(x) \varphi(x) - \frac{v^2}{2} \right)^2
\end{align*}
Since $\varphi^\dagger \varphi$ remains invariant under the gauge transformation, we have:
\begin{align*}
    \left( \varphi^\dagger(x) \varphi(x) - \frac{v^2}{2} \right) &\to \left( \varphi^\dagger(x) \varphi(x) - \frac{v^2}{2} \right)
\end{align*}
Therefore, the entire potential term transforms as:
\begin{align*}
    \left( \varphi^\dagger(x) \varphi(x) - \frac{v^2}{2} \right)^2 &\to \left( \varphi^\dagger(x) \varphi(x) - \frac{v^2}{2} \right)^2
\end{align*}
This shows that the potential term remains unchanged under the gauge transformation. Finally, the Lagrangian term becomes:
\begin{align*}
    L_{\text{potential}} &= - \frac{\lambda}{4} \left( \varphi^\dagger(x) \varphi(x) - \frac{v^2}{2} \right)^2
\end{align*}
Thus, we conclude that the scalar potential term is gauge-invariant under a local $U(1)$ transformation.
\item [(iv)] \textbf{Yukawa Interactions}

We start with the Yukawa terms in the Lagrangian:
\begin{align}
    L_{\text{Yukawa}} &= - \left( \Gamma_2 \bar{\psi} \varphi P_R \psi_2 + \Gamma_1 \bar{\psi} \varphi^c P_R \psi_1 + \text{h.c.} \right)\label{eq:yukawa-lagrangian-def}
\end{align}
The fermion fields $\psi$, $\psi_1$, and $\psi_2$ transform under a local $U(1)$ gauge transformation as:
\begin{align*}
    \psi(x) &\to e^{\displaystyle i \alpha(x)} \psi(x) \\
    \bar{\psi}(x) &\to e^{\displaystyle -i \alpha(x)} \bar{\psi}(x) \\
    \psi_1(x) &\to e^{\displaystyle i \alpha(x)} \psi_1(x) \\
    \psi_2(x) &\to e^{\displaystyle i \alpha(x)} \psi_2(x)
\end{align*}
The scalar field $\varphi(x)$ transforms as:
\begin{align*}
    \varphi(x) &\to e^{\displaystyle i \alpha(x)} \varphi(x)
\end{align*}
The conjugate scalar field $\varphi^c(x)$ transforms as:
\begin{align*}
    \varphi^c(x) &\to e^{\displaystyle -i \alpha(x)} \varphi^c(x)
\end{align*}
Now, we analyze how the bilinear terms involving the fermions and scalar fields transform.

For the first term $\bar{\psi} \varphi P_R \psi_2$:
\begin{align*}
    \bar{\psi}(x) \varphi(x) P_R \psi_2(x) &\to e^{\displaystyle -i \alpha(x)} \bar{\psi}(x) e^{i \alpha(x)} \varphi(x) P_R e^{\displaystyle i \alpha(x)} \psi_2(x)
\end{align*}
The exponentials cancel out, leaving:
\begin{align*}
    \bar{\psi}(x) \varphi(x) P_R \psi_2(x) &\to \bar{\psi}(x) \varphi(x) P_R \psi_2(x)
\end{align*}
Thus, the first term is invariant under the gauge transformation.

For the second term $\bar{\psi} \varphi^c P_R \psi_1$:
\begin{align*}
    \bar{\psi}(x) \varphi^c(x) P_R \psi_1(x) &\to e^{\displaystyle -i \alpha(x)} \bar{\psi}(x) e^{\displaystyle -i \alpha(x)} \varphi^c(x) P_R e^{\displaystyle i \alpha(x)} \psi_1(x)
\end{align*}
Again, the exponentials cancel out, giving:
\begin{align*}
    \bar{\psi}(x) \varphi^c(x) P_R \psi_1(x) &\to \bar{\psi}(x) \varphi^c(x) P_R \psi_1(x)
\end{align*}
Thus, the second term is also invariant under the gauge transformation.
\end{enumerate}
A bare mass term for the gauge bosons or fermions would take the form $m \bar{\psi} \psi$ or $m W_\mu^a W^{a\mu}$. Such terms are not allowed because they are not gauge invariant.

For example, a bare mass term for the fermions:
\begin{align*}
    L_{\text{bare fermion mass}} &\sim \bar{\psi} \psi
\end{align*}
is not invariant under a gauge transformation $\psi \to U(x) \psi$, since it would transform as:
\begin{align*}
    m \bar{\psi} \psi \to m \bar{\psi} U^\dagger(x) U(x) \psi = m \bar{\psi} \psi
\end{align*}
which is not invariant unless $U(x) = 1$ everywhere, implying no gauge symmetry.

Similarly, a bare mass term for the gauge bosons:
\begin{align*}
    L_{\text{bare gauge boson mass}} &\sim W_\mu^a W^{a\mu}
\end{align*}
is also not gauge invariant because the gauge field transforms non-trivially under $\text{SU(2)}$:
\begin{align*}
    W_\mu^a \to W_\mu^{a^\prime } = W_\mu^a + \frac{1}{g} \partial_\mu \alpha^a + \epsilon^{abc} \alpha^b W_\mu^c
\end{align*}
Hence, the mass term is not invariant under this transformation.

A bare mass term for the scalar field $\varphi$ would take the form:
\begin{align}
    L_{\text{bare scalar mass}} &\sim - m_\varphi^2 \varphi^\dagger \varphi
\end{align}
This term represents a mass for the scalar field $\varphi$ without any interaction or symmetry-breaking mechanism.

Under a local $\text{SU(2)} \otimes \text{U(1)}$ gauge transformation, the scalar field $\varphi$ transforms as:
\begin{align*}
    \varphi(x) &\to U(x) \varphi(x) = e^{\displaystyle i \alpha^a(x) \frac{\tau^a}{2}} e^{\displaystyle i \beta(x) \frac{Y}{2}} \varphi(x)
\end{align*}
where $U(x)$ is the local gauge transformation matrix, $\alpha^a(x)$ is the parameter for $\text{SU(2)}$, and $\beta(x)$ is the parameter for $\text{U(1)}$.

The scalar mass term transforms as:
\begin{align*}
    \varphi^\dagger \varphi &\to \varphi^\dagger U^\dagger(x) U(x) \varphi = \varphi^\dagger \varphi
\end{align*}
Since $U(x)$ is unitary, $U^\dagger(x) U(x) = 1$, and thus $\varphi^\dagger \varphi$ is invariant under the gauge transformation.

Thus, no bare mass terms are allowed by gauge invariance.

We now try to see the effects of spontaneous symmetry breaking.

The scalar doublet $\varphi$ is a complex field transforming under $\text{SU(2)} \otimes \text{U(1)}$. It acquires a non-zero vacuum expectation value (VEV) due to spontaneous symmetry breaking:
\begin{align}
    \varphi &= \frac{1}{\sqrt{2}} \begin{pmatrix} 0 \\ v + h(x) \end{pmatrix}
\end{align}
where $v$ is the VEV and $h(x)$ represents the physical Higgs field. The VEV breaks the gauge symmetry $\text{SU(2)} \otimes \text{U(1)}$ to the electromagnetic subgroup $\text{U(1)}_{\text{em}}$.

The covariant derivative for the scalar field $\varphi$ is given by:
\begin{align}
    D_\mu \varphi &= \left( \partial_\mu - i g \frac{\tau^a}{2} W_\mu^a - i g^\prime \frac{Y}{2} B_\mu \right) \varphi
\end{align}
where $g$ is the $\text{SU(2)}$ gauge coupling, $g^\prime $ is the $\text{U(1)}$ gauge coupling, $\tau^a$ are the Pauli matrices, and $Y$ is the hypercharge of the scalar doublet ($Y = 1$). 

Now, simplify each part of the expression. For the $\partial_\mu$ term:
\begin{align}
    \partial_\mu \varphi &= \frac{1}{\sqrt{2}} \begin{pmatrix} 0 \\\\ \partial_\mu (v + h(x)) \end{pmatrix}
\end{align}
For the $\text{SU(2)}$ gauge term:
\begin{align}
    - i \frac{g}{2} \begin{pmatrix} W_\mu^3 & W_\mu^1 - i W_\mu^2 \\\\ W_\mu^1 + i W_\mu^2 & -W_\mu^3 \end{pmatrix} \begin{pmatrix} 0 \\\\ \displaystyle\frac{v + h(x)}{\sqrt{2}} \end{pmatrix} &= - i \frac{g}{2} \frac{v + h(x)}{\sqrt{2}} \begin{pmatrix} W_\mu^1 + i W_\mu^2 \\\\ -W_\mu^3 \end{pmatrix}
\end{align}
For the $\text{U(1)}$ gauge term:
\begin{align}
    - i \frac{g^\prime }{2} B_\mu \begin{pmatrix} 0 \\\\ \displaystyle\frac{v + h(x)}{\sqrt{2}} \end{pmatrix} &= - i \frac{g^\prime }{2} \frac{v + h(x)}{\sqrt{2}} \begin{pmatrix} 0 \\\\ B_\mu \end{pmatrix}
\end{align}
Thus, the covariant derivative acting on $\varphi$ is:
\begin{align}
    D_\mu \varphi &= \frac{1}{\sqrt{2}} \begin{pmatrix} 0 \\ \\ \displaystyle\partial_\mu (v + h(x)) \end{pmatrix} - i \frac{v + h(x)}{\sqrt{2}} \begin{pmatrix} \displaystyle\frac{g}{2} (W_\mu^1 + i W_\mu^2) \\ \\ \displaystyle\frac{g}{2} W_\mu^3 + \frac{g^\prime }{2} B_\mu \end{pmatrix}
\end{align}
The Hermitian conjugate of $D_\mu \varphi$ is given by:
\begin{align}
    (D_\mu \varphi)^\dagger &= \frac{1}{\sqrt{2}} \left( 0, \partial_\mu (v + h(x)) \right)^\dagger + i \frac{v + h(x)}{\sqrt{2}} \left( \frac{g}{2} (W_\mu^1 - i W_\mu^2), \frac{g}{2} W_\mu^3 + \frac{g^\prime }{2} B_\mu \right)
\end{align}
We now compute the product $(D_\mu \varphi)^\dagger (D^\mu \varphi)$:
\begin{align}
    (D_\mu \varphi)^\dagger (D^\mu \varphi) &= \frac{v^2}{2} \left[ \frac{g^2}{4} \left( W_\mu^1 W^{1\mu} + W_\mu^2 W^{2\mu} \right) + \frac{1}{4} (g W_\mu^3 + g^\prime B_\mu)^2 \right]
\end{align}
This expression can be simplified further. The mass terms for the $W^\pm$ bosons arise from the terms involving $W_\mu^1$ and $W_\mu^2$. We introduce the charged $W^\pm$ fields as:
\begin{align}
    W_\mu^\pm &= \frac{1}{\sqrt{2}} (W_\mu^1 \mp i W_\mu^2)
\end{align}
This gives:
\begin{align*}
    W_\mu^1 W^{1\mu} + W_\mu^2 W^{2\mu} &= 2 W^{+\dagger}_\mu W^{\mu+}
\end{align*}
Thus, the $W$ boson mass term becomes:
\begin{align}
    \frac{g^2 v^2}{4} W^{+\dagger}_\mu W^{\mu+} &= m_W^2 W^{+\dagger}_\mu W^{\mu+}
\end{align}
where the $W$ boson mass is:
\begin{align}
    m_W &= \frac{gv}{2}
\end{align}
The mass term for the $Z$ boson arises from the combination of $W_\mu^3$ and $B_\mu$:
\begin{align}
    \frac{1}{4} (g W_\mu^3 - g^\prime B_\mu)^2 &= \frac{1}{4} (g^2 W_\mu^3 W^{3\mu} - 2 g g^\prime W_\mu^3 B^\mu + g^\prime ^2 B_\mu B^\mu)
\end{align}
Define the $Z_\mu$ and $A_\mu$ fields as:
\begin{align}
    Z_\mu &= \cos(\theta_W) W_\mu^3 - \sin(\theta_W) B_\mu \label{eq:Z-mu-def}\\
    A_\mu &= \sin(\theta_W) W_\mu^3 + \cos(\theta_W) B_\mu \label{eq:A-mu-def}
\end{align}
where the Weinberg angle $\theta_W$ is defined as:
\begin{align}
    \cos(\theta_W) &= \frac{g}{\sqrt{g^2 + g^\prime ^2}} \label{eq;cos-weinberg-def}\\
    \sin(\theta_W) &= \frac{g^\prime }{\sqrt{g^2 + g^\prime ^2}} \label{eq;sin-weinberg-def}
\end{align}
This leads to the mass term for the $Z$ boson:
\begin{align}
    \frac{v^2}{2} \frac{g^2 + g^\prime ^2}{4} Z_\mu Z^\mu &= \frac{1}{2} m_Z^2 Z_\mu Z^\mu
\end{align}
where the $Z$ boson mass is:
\begin{align}
    m_Z &= \frac{\sqrt{g^2 + g^\prime ^2}}{2} v
\end{align}
The gauge boson mass terms are:
\begin{align}
    m_W^2 W^{+\dagger}_\mu W^{\mu+} + \frac{1}{2} m_Z^2 Z_\mu Z^\mu
\end{align}
where the masses of the $W$ and $Z$ bosons are:
\begin{align}
    m_W &= \frac{gv}{2} \label{eq:m_W-def}\\
    m_Z &= \frac{\sqrt{g^2 + g^\prime ^2}}{2} v\label{eq:m_Z-def}
\end{align}
The ratio $\displaystyle \frac{m_Z}{m_W}$ is, using (\ref{eq:m_W-def}), (\ref{eq:m_Z-def}) and (\ref{eq;cos-weinberg-def}):
\begin{align}
    \frac{m_Z}{m_W} &= \frac{\displaystyle\frac{\sqrt{g^2 + g^\prime ^2}}{2} v}{\displaystyle\frac{gv}{2}} = \frac{\sqrt{g^2 + g^\prime ^2}}{g} = \frac{1}{\cos(\theta_W)}.\qquad(\textbf{Answer})
\end{align}
We substitute the vacuum expectation value (VEV) of the scalar field $\varphi$, and its charge conjugates $\varphi^c$ as:
\begin{align}
    \varphi &= \frac{1}{\sqrt{2}} \begin{pmatrix} 0 \\\\ v + h(x) \end{pmatrix}, \quad \varphi^c = \frac{1}{\sqrt{2}} \begin{pmatrix} v + h(x) \\\\ 0 \end{pmatrix}
\end{align}
Substituting the VEV $\varphi$ and $\varphi^c$ into the Yukawa Lagrangian (\ref{eq:yukawa-lagrangian-def}):
\begin{align}
    L_{\text{Yukawa}} &= - \left( \Gamma_2 \bar{\psi} \frac{1}{\sqrt{2}} \begin{pmatrix} 0 \\\\ v + h(x) \end{pmatrix} P_R \psi_2 + \Gamma_1 \bar{\psi} \frac{1}{\sqrt{2}} \begin{pmatrix} v + h(x) \\\\ 0 \end{pmatrix} P_R \psi_1 + \text{h.c.} \right)
\end{align}
For the first term involving $\varphi$:
\begin{align}
    \Gamma_2 \bar{\psi} \varphi P_R \psi_2 &= \Gamma_2 \bar{\psi} \frac{1}{\sqrt{2}} \begin{pmatrix} 0 \\\\ v + h(x) \end{pmatrix} P_R \psi_2 = \frac{\Gamma_2}{\sqrt{2}} \bar{\psi} (v + h(x)) P_R \psi_2
\end{align}
For the second term involving $\varphi^c$:
\begin{align}
    \Gamma_1 \bar{\psi} \varphi^c P_R \psi_1 &= \Gamma_1 \bar{\psi} \frac{1}{\sqrt{2}} \begin{pmatrix} v + h(x) \\\\ 0 \end{pmatrix} P_R \psi_1 = \frac{\Gamma_1}{\sqrt{2}} \bar{\psi} (v + h(x)) P_R \psi_1
\end{align}
Now substitute the expanded terms back into the Lagrangian:
\begin{align}
    L_{\text{Yukawa}} &= - \left( \frac{\Gamma_2}{\sqrt{2}} \bar{\psi} (v + h(x)) P_R \psi_2 + \frac{\Gamma_1}{\sqrt{2}} \bar{\psi} (v + h(x)) P_R \psi_1 + \text{h.c.} \right)
\end{align}
We now separate the terms involving the VEV $v$ and the Higgs field fluctuation $h(x)$:
\begin{align}
    L_{\text{Yukawa}} &= - \frac{v}{\sqrt{2}} \left( \Gamma_2 \bar{\psi} P_R \psi_2 + \Gamma_1 \bar{\psi} P_R \psi_1 \right) - \frac{h(x)}{\sqrt{2}} \left( \Gamma_2 \bar{\psi} P_R \psi_2 + \Gamma_1 \bar{\psi} P_R \psi_1 \right) + \text{h.c.}
\end{align}
From the terms involving the VEV $v$, we identify the fermion masses:
\begin{align}
    m_{\psi_2} &= \frac{\Gamma_2 v}{\sqrt{2}} \\
    m_{\psi_1} &= \frac{\Gamma_1 v}{\sqrt{2}}
\end{align}
Thus, the full Yukawa Lagrangian (\ref{eq:yukawa-lagrangian-def}), including both the mass terms and the interaction with the Higgs field $h(x)$, becomes:
\begin{align}
    L_{\text{Yukawa}} &= - m_{\psi_2} \bar{\psi} P_R \psi_2 - m_{\psi_1} \bar{\psi} P_R \psi_1 - \frac{h(x)}{\sqrt{2}} \left( \Gamma_2 \bar{\psi} P_R \psi_2 + \Gamma_1 \bar{\psi} P_R \psi_1 \right) + \text{h.c.}
\end{align}

The terms $m_{\psi_2}$ and $m_{\psi_1}$ are the masses for the fermions $\psi_2$ and $\psi_1$, respectively, while the interaction terms involving $h(x)$ describe the Yukawa interaction between the fermions and the Higgs field.
\bigskip\bigskip\hline\hline\bigskip